\section{Applying the Model: Evaluating the Effectiveness of Visualization Descriptions}
\label{section:5}

The foregoing conceptual model provides a means of making structured comparisons between different levels of semantic content and reader groups.
To demonstrate how it can be applied to evaluate the effectiveness of visualization descriptions (i.e., whether or not they effectively convey meaningful information, and for whom), we conducted a mixed-methods evaluation in which {30} blind and {90} sighted readers first ranked the usefulness of descriptions authored at varying levels of semantic content, and then completed an open-ended questionnaire.

%%%%%%%%%%%%%%%%%%%%%%%%%%%%%%%%%%%%%%%%%%%%%%%%%%%%%%%%%%%%%%%%%
% Evaluation Design
%%%%%%%%%%%%%%%%%%%%%%%%%%%%%%%%%%%%%%%%%%%%%%%%%%%%%%%%%%%%%%%%%
\subsection{Evaluation Design}
\label{section:study2design}

We selected 15 visualizations for the evaluation, curated to be representative of the categories from our prior survey~\sectionlabel{section:3}.
Specifically, we selected 5 visualizations for each of the three dimensions: \emph{type} (bar, line, scatter), \emph{topic} (academic, business, journalism), and \emph{difficulty} (easy, medium, hard).
For every visualization, participants were asked to rank the usefulness of 4 different descriptions, each corresponding to one level of semantic content, presented unlabeled and in random order.
We piloted this rank-choice interface with 10 sighted readers recruited via \emph{Prolific} and 1 blind reader, a non-academic collaborator proficient with Apple's VoiceOver screen reader. 
Based on this pilot, we rewrote the study instructions to be more intelligible to both groups of readers, added an introductory example task to the evaluation, and improved the screen reader accessibility of our interface (e.g., by reordering nested \textsc{dom} elements to be more intuitively traversed by screen reader).

In addition to curating a representative set of visualizations, we also curated descriptions representative of each level of semantic content.
Participant-authored descriptions from our prior survey often did not contain content from all 4 levels or, if they did, this content was interleaved in a way that was not cleanly-separable for the purpose of a ranking task (Fig.~\ref{fig:fingerprint}).
Thus, for this evaluation, we curated and collated sentences from multiple participant-authored descriptions to create exemplar descriptions, such that each text chunk contained \emph{only} content belonging to a single semantic content level.
Table~\ref{table:visualizations}.C shows one such exemplar description, whereas Table~\ref{table:visualizations}.A and B show the original un-collated descriptions.
For each ranking task, readers were presented with a brief piece of contextualizing text, such as the following.
\begin{displayquote}
    % \small
    \vspace{-1mm}
    \say{\emph{Suppose that you are reading an academic paper about how life expectancy differs for people of different genders from countries with different levels of income. You encounter the following visualization.} [Table~\ref{table:visualizations}.C]
    \emph{Which content do you think would be most useful to include in a textual description of this visualization?}}
    \vspace{-1mm}
\end{displayquote}
Additionally, blind readers were presented with a brief text noting that the hypothetically-encountered visualization was inaccessible via screen reader technology.
In contrast to prior work, which has evaluated chart descriptions in terms of ``efficiency,'' ``informativeness,'' and ``clarity''~\cite{gould_effective_2008, obeid_chart--text_2020}, we intentionally left the definition of ``useful'' open to the reader’s interpretation. 
We hypothesize that ``useful'' descriptions may not be necessarily efficient (i.e., they may require lengthy explanation or background context), and that both informativeness and clarity are constituents of usefulness.
In short, ranking ``usefulness'' affords a holistic evaluation metric.
Participants assigned usefulness rankings to each of the 4 descriptions by selecting corresponding radio buttons, {labeled 1 (least useful) to 4 (most useful)}.
In addition to these 4 descriptions, we included a 5th choice as an ``attention check'': a sentence whose content was entirely irrelevant to the chart to ensure participants were reading each description prior to ranking them.
If a participant did not rank the attention check as least useful, we filtered out their response from our final analysis.
{We include the evaluation interfaces and questions with the Supplemental Material.}

%%%%%%%%%%%%%%%%%%%%%%%%%%%%%%%%%%%%%%%%%%%%%%%%%%%%%%%%%%%%%%%%%
% Participants
%%%%%%%%%%%%%%%%%%%%%%%%%%%%%%%%%%%%%%%%%%%%%%%%%%%%%%%%%%%%%%%%%
\subsection{Participants}

Participants consisted of two reader groups: {90} sighted readers recruited through the \emph{Prolific} platform, and {30} blind readers recruited through our friends in the blind community and through a call for participation sent out via Twitter (n.b., in accessibility research, it is common to compare blind and sighted readers recruited through these means~\cite{bigham_effects_2017}).

\subsubsection{Participant Recruitment}

For sighted readers qualifications for participation included English language proficiency and no color vision deficiency, and blind readers were expected to be proficient with a screen reader, such as \ac{JAWS}, \ac{NVDA}, or Apple's VoiceOver.
Sighted readers were compensated at a rate of \$10-12 per hour, for an approximately 20-minute task. 
Blind readers were compensated at a rate of \$50 per hour, for an approximately 1-hour task. 
This difference in task duration was for two reasons. 
First, participants recruited through \emph{Prolific} are usually not accustomed to completing lengthy tasks\,---\,our prior surveys and pilots suggested that these participants might contribute low-quality responses on ``click-through'' tasks if the task duration exceeded 15--20 minutes\,---\,and thus we asked each participant to rank only 5 of the 15 visualizations at a time.
Second, given the difficulty of recruiting blind readers proficient with screen readers, we asked each blind participant to rank all 15 visualizations, and compensated them at a rate commensurate with their difficult-to-find expertise~\cite{lundgard_sociotechnical_2019}.
In this way, we recruited sufficient numbers of readers to ensure that each of the 15 visualization ranking tasks would be completed by 30 participants from both reader groups.

\subsubsection{Participant Demographics}
Among the {30} blind participants, {53}\% (n={16}) reported their gender as male, {36}\% (n={11}) as female, and {3} participants ``preferred not to say.''
The most common highest level of education attained was a Bachelor's degree ({60}\%, n=18), and most readers were between 20\,--\,40 years old ({66}\%, n=20).
The screen reader technology readers used to complete the study was evenly balanced: VoiceOver (n=10), \ac{JAWS} (n=10), \ac{NVDA} (n=9), and ``other'' (n=1).
Among the {90} sighted participants, {69}\% reported their gender as male (n={62}) and {31}\% as female (n={28}).
The most common highest level of education attained was a high school diploma ({42}\%, n=38) followed by a Bachelor's degree ({40}\%, n=36), and most sighted readers were between 20\,--\,30 years old ({64}\%, n=58).

On a 7-point Likert scale [1=strongly disagree, 7=strongly agree], blind participants reported having ``a good understanding of data visualization concepts'' ($\mu=6.3$, $\sigma=1.03$) as well as ``a good understanding of statistical concepts and terminology'' ($\mu=5.90$, $\sigma=1.01$). Sighted participants reported similar levels of understanding: ($\mu=6.7$, $\sigma=0.73$) and ($\mu=5.67$, $\sigma=1.06$), respectively.
Sighted participants also considered themselves to be ``proficient at reading data visualizations'' ($\mu=5.97$, $\sigma=0.89$) and were able to ``read and understand all of the visualizations presented in this study'' ($\mu=6.44$, $\sigma=0.71$).

% Most readers were between ages 30-39 (n=11), followed by ages 40-49 (n=7), ages 20-29 (n=5), and ages 50-59 (n=1).
% The most common reported occupation was Assistive Technology Specialist or Engineer (n=6), with other occupations including Nonprofit Program Director, Applied Linguist, Artist, and Lawyer.
% The screen reader technology participants used to complete the study was evenly balanced: \ac{JAWS} (n=9), \ac{NVDA} (n=8), VoiceOver (n=8), and "Other" (n=1). 

% Most readers were between ages 30-39 (n=11), followed by ages 40-49 (n=7), ages 20-29 (n=5), and ages 50-59 (n=1).
% Reported occupations included Computer Scientist, Graphic Designer, Taxi Driver, German Philologist, Information Technologist, among many others.

%%%%%%%%%%%%%%%%%%%%%%%%%%%%%%%%%%%%%%%%%%%%%%%%%%%%%%%%%%%%%%%%%
% Evaluation Table
%%%%%%%%%%%%%%%%%%%%%%%%%%%%%%%%%%%%%%%%%%%%%%%%%%%%%%%%%%%%%%%%%
% Evaluation Table Vertical
%%%%%%%%%%%%%%%%%%%%%%%%%%%%%%%%%%%%%%%%%%%%%%%%%%%%%%%%%%%%%%%%%
\bgroup
\begin{table}
\centering
\caption{
(Upper) Rankings [1=least useful, 4=most useful] of semantic content at each level of the model, for blind and sighted readers.
The scale encodes the number of times a given level was assigned a given rank by a reader.
Dotted contour lines delineate Regions with a threshold equal to $\mu + \frac{\sigma}{2}$, each labeled with a capital letter  A\,--\,F.
%in its upper-right-hand corner.
(Lower)
% Pair-wise ranking differences between levels. 
Shaded cells indicate significant ranking differences pair-wise between levels.
}

\def\arraystretch{1.5}
\setlength{\tabcolsep}{2pt}
\smaller

\vspace{-1mm}
%---------------------------------------------------------------
% Heatmaps
%---------------------------------------------------------------
\begin{tabular}{cc}
\hline
\hlc[DirectQuoteA]{\normalsize\textsc{blind readers}} & \hlc[DirectQuoteB]{\normalsize\textsc{sighted readers}} \\
\hline
\includegraphics[width=0.22\textwidth]{figures/final_blind_rankings_numeric_anon.png}
\vspace{-1mm}
&
\includegraphics[width=0.22\textwidth]{figures/final_sighted_rankings_numeric_anon.png}
\\
% \hline
\end{tabular}
% ---------------------------------------------------------------
% ---------------------------------------------------------------
\vspace{6mm}
%---------------------------------------------------------------
% Significance
%---------------------------------------------------------------
\begin{tabular}{c|c|c|c|c|c|c}
\hline
\textsc{levels} & $1 \times 2$ & $1 \times 3$ & $1 \times 4$ & $2 \times 3$ & $2 \times 4$ & $3 \times 4$ \\ \hline
\hlc[DirectQuoteA]{\textsc{blind}} & \cellcolor{yellow!25}$p<0.001$ & \cellcolor{yellow!25}$p<0.001$ & $p<0.321$ & $p<0.148$ & \cellcolor{yellow!25}$p<0.001$ & \cellcolor{yellow!25}$p<0.001$ \\ \hline
\hlc[DirectQuoteB]{\textsc{sighted}} & \cellcolor{yellow!25}$p<0.001$ & \cellcolor{yellow!25}$p<0.001$ & \cellcolor{yellow!25}$p<0.001$ & \cellcolor{yellow!25}$p<0.001$ & \cellcolor{yellow!25}$p<0.001$ & $p<0.059$ \\ \hline
\end{tabular}
%---------------------------------------------------------------
%---------------------------------------------------------------
\label{table:evaluation}
\vspace{-11mm} %FIX ME? Camera-ready
\end{table}
\egroup
%%%%%%%%%%%%%%%%%%%%%%%%%%%%%%%%%%%%%%%%%%%%%%%%%%%%%%%%%%%%%%%%%

% %%%%%%%%%%%%%%%%%%%%%%%%%%%%%%%%%%%%%%%%%%%%%%%%%%%%%%%%%%%%%%%%%
% % Evaluation Table (Heatmaps Only)
% %%%%%%%%%%%%%%%%%%%%%%%%%%%%%%%%%%%%%%%%%%%%%%%%%%%%%%%%%%%%%%%%%
% \bgroup
% \begin{table}[]
% \centering
% \caption{
% Rankings [1=least useful, 4=most useful] of semantic content at each level of the model, for blind and sighted readers.
% The scale encodes the number of times a level was assigned a rank by a reader.
% The dotted contour lines delineate regions with a threshold equal to $\mu + \frac{\sigma}{2}$.
% Each region is labeled with a capital letter  A\,--\,F in its upper-right-hand corner.
% }
% \def\arraystretch{1.5}
% \vspace{-1mm}
% \small

% %---------------------------------------------------------------
% % Heatmaps
% %---------------------------------------------------------------
% \begin{tabular}{cc}
% \hline
% \hlc[DirectQuoteA]{\textsc{blind readers}} & \hlc[DirectQuoteB]{\textsc{sighted readers}} \\
% \hline
% \includegraphics[width=0.21\textwidth]{figures/final_blind_rankings_numeric_anon.png}
% \vspace{-1mm}
% &
% \includegraphics[width=0.21\textwidth]{figures/final_sighted_rankings_numeric_anon.png}
% \vspace{-1mm}
% \\
% \hline
% \end{tabular}
% %---------------------------------------------------------------
% %---------------------------------------------------------------

% \label{table:evaluation}
% % \vspace{-5mm}
% \end{table}
% \egroup
% %%%%%%%%%%%%%%%%%%%%%%%%%%%%%%%%%%%%%%%%%%%%%%%%%%%%%%%%%%%%%%%%%

% %%%%%%%%%%%%%%%%%%%%%%%%%%%%%%%%%%%%%%%%%%%%%%%%%%%%%%%%%%%%%%%%%
% % Evaluation Table Horizontal
% %%%%%%%%%%%%%%%%%%%%%%%%%%%%%%%%%%%%%%%%%%%%%%%%%%%%%%%%%%%%%%%%%
% \bgroup
% \begin{table*}[]
% \centering
% \caption{
% (Left) Pair-wise ranking differences between levels, for blind and sighted readers. Highlighted cells show significant ranking differences.
% (Right) Rankings [1=least useful, 4=most useful] of semantic content at each level of the model, for blind and sighted readers.
% The scale encodes the number of times a given level was assigned a given rank by a reader.
% The dotted contour lines delineate regions with a threshold equal to $\mu + \frac{\sigma}{2}$.
% Each region is labeled with a capital letter  A\,--\,F in its upper-right-hand corner.
% }
% \def\arraystretch{1.5}
% \vspace{-1mm}
% \small
% %---------------------------------------------------------------
% % Significance
% %---------------------------------------------------------------
% \begin{tabular}{ccc}
% \hline
% % $(level\times level)$
% \textsc{levels} & \hlc[DirectQuoteA]{\textsc{blind}} & \hlc[DirectQuoteB]{\textsc{sighted}} \\ \hline
% $1 \times 2$ & \cellcolor{yellow!25}$p<0.001$ & \cellcolor{yellow!25}$p<0.001$ \\ \hline
% $1 \times 3$ & \cellcolor{yellow!25}$p<0.001$ & \cellcolor{yellow!25}$p<0.001$ \\ \hline
% $1 \times 4$ & $p<0.321$ & \cellcolor{yellow!25}$p<0.001$ \\ \hline
% $2 \times 3$ & $p<0.148$ & \cellcolor{yellow!25}$p<0.001$ \\ \hline
% $2 \times 4$ & \cellcolor{yellow!25}$p<0.001$ & \cellcolor{yellow!25}$p<0.001$ \\ \hline
% $3 \times 4$ & \cellcolor{yellow!25}$p<0.001$ & $p<0.059$ \\ \hline
% \end{tabular}
% %---------------------------------------------------------------
% %---------------------------------------------------------------
% \hspace{5mm}
% %---------------------------------------------------------------
% % Heatmaps
% %---------------------------------------------------------------
% \begin{tabular}{cc}
% \hline
% \hlc[DirectQuoteA]{\textsc{blind readers}} & \hlc[DirectQuoteB]{\textsc{sighted readers}} \\
% \hline
% \includegraphics[width=0.25\textwidth]{figures/final_blind_rankings_numeric_anon.png}
% \vspace{-1mm}
% &
% \includegraphics[width=0.25\textwidth]{figures/final_sighted_rankings_numeric_anon.png}
% \vspace{-1mm}
% \\
% \hline
% \end{tabular}
% %---------------------------------------------------------------
% %---------------------------------------------------------------

% \label{table:evaluation}
% \vspace{-5mm}
% \end{table*}
% \egroup
% %%%%%%%%%%%%%%%%%%%%%%%%%%%%%%%%%%%%%%%%%%%%%%%%%%%%%%%%%%%%%%%%%

% %%%%%%%%%%%%%%%%%%%%%%%%%%%%%%%%%%%%%%%%%%%%%%%%%%%%%%%%%%%%%%%%%
% % Evaluation Table Vertical Composite
% %%%%%%%%%%%%%%%%%%%%%%%%%%%%%%%%%%%%%%%%%%%%%%%%%%%%%%%%%%%%%%%%%
% \bgroup
% \begin{table*}[]

% \caption{Rankings (1=least useful, 4=most useful) of descriptive content at each level of the hierarchy, for reader groups (blind or sighted).
% The scale encodes the number of times a given description level was assigned a given rank by readers.
% The dotted contour lines delineate regions of the heatmap with a threshold=140.
% Each region is labeled with a captial letter in its upper-right-hand corner.
% \hl{TODO: Revise.}
% }

% \centering
% \def\arraystretch{1.5}

% \begin{tabular}{cc}
% \multirow{4}{*}{
% \includegraphics[width=0.61\linewidth]{figures/visual_fingerprint_comp.png}
% } 
% & \cellcolor{DirectQuoteA}{\textsc{blind readers}} \\
% &\includegraphics[width=0.3\textwidth]{figures/final_blind_rankings_numeric_anon.png}
% \\
% &\cellcolor{DirectQuoteB}{\textsc{sighted readers}} \\
% &\includegraphics[width=0.3\textwidth]{figures/final_sighted_rankings_numeric_anon.png}
% \\
% \end{tabular}

% \label{table:evaluation}
% \vspace{-5mm}
% \end{table*}
% \egroup
% %%%%%%%%%%%%%%%%%%%%%%%%%%%%%%%%%%%%%%%%%%%%%%%%%%%%%%%%%%%%%%%%%

%%%%%%%%%%%%%%%%%%%%%%%%%%%%%%%%%%%%%%%%%%%%%%%%%%%%%%%%%%%%%%%%%

%%%%%%%%%%%%%%%%%%%%%%%%%%%%%%%%%%%%%%%%%%%%%%%%%%%%%%%%%%%%%%%%%
% Quantitative Results
%%%%%%%%%%%%%%%%%%%%%%%%%%%%%%%%%%%%%%%%%%%%%%%%%%%%%%%%%%%%%%%%%
\subsection{Quantitative Results}

% \noindent
% \textbf{\emph{Ranked Preferences}}

Quantitative results for the individual rankings ({1,800} per blind and sighted reader groups) are summarized by the heatmaps in Table~\ref{table:evaluation} (Upper), which aggregate the number of times a given content level was assigned a certain rank.
Dotted lines in both blind and sighted heatmaps delineate regions exceeding a threshold\,---\,calculated by taking the mean plus half a standard deviation ($\mu + \frac{\sigma}{2}$) resulting in a value of 139 and 136, respectively\,---\,and are labeled with a capital letter A\,--\,F.

These results exhibit significant differences between reader groups.
For both reader groups, using Friedman’s Test (a non-parametric multi-comparison test for rank-order data) the p-value is $p<0.001$, so we reject the null hypothesis that the mean rank is the same for all four semantic content levels~\cite{garcia_advanced_2010}.
Additionally, in Table~\ref{table:evaluation} (Lower), we find significant ranking differences when making pair-wise comparisons between levels, via Nemenyi’s test (a post-hoc test commonly coupled with Friedman's to make pair-wise comparisons).
There appears to be strong agreement among sighted readers that higher levels of semantic content are more useful: Levels 3 and 4 are found to be most useful (Region~\ref{table:evaluation}.F), while Levels 1 and 2 are least useful (Regions~\ref{table:evaluation}.D and ~\ref{table:evaluation}.E).
Blind readers agree with each other to a lesser extent, but strong trends are nevertheless apparent.
In particular, blind readers rank content and Levels 2 and 3 as most useful (Region~\ref{table:evaluation}.C), and semantic content at Levels 1 and 4 as least useful (Regions~\ref{table:evaluation}.A and ~\ref{table:evaluation}.B).

When faceting these rankings by visualization type, topic, or difficulty we did not observe any significant differences, suggesting that both reader groups rank semantic content levels consistently, regardless of how the chart itself may vary.
Noteworthy for both reader groups, the distribution of rankings for Level 1 is bimodal\,—--\,the only level to exhibit this property.
While a {vast majority} of both blind and sighted readers rank Level 1 content as least useful, this level is ranked ``most useful'' in {101} and {87} instances by blind and sighted readers, respectively. This suggests that both reader groups have a more complicated perspective toward descriptions of a chart's elemental and encoded properties; a finding we explore further by analyzing qualitative data.

%%%%%%%%%%%%%%%%%%%%%%%%%%%%%%%%%%%%%%%%%%%%%%%%%%%%%%%%%%%%%%%%%
% Qualitative Results
%%%%%%%%%%%%%%%%%%%%%%%%%%%%%%%%%%%%%%%%%%%%%%%%%%%%%%%%%%%%%%%%%
\subsection{Qualitative Results}
\label{section:qualResults}

In a questionnaire, we asked readers to use a 7-point Likert scale [1=strongly disagree, 7=strongly agree] to rate their agreement with a set of statements about their experience with visualizations.
We also asked them to offer open-ended feedback about which semantic content they found to be most useful and why.
Here, we summarize the key trends that emerged from these two different forms of feedback, from both \hlc[DirectQuoteA]{blind readers} (BR) and \hlc[DirectQuoteB]{sighted readers} (SR).

%----------------------------------------------------------------
\subsubsection{Descriptions Are Important to Both Reader Groups}
\label{section:qualResultsA}
%----------------------------------------------------------------
% Responding to the statement \say{Providing textual descriptions of data visualizations is important to me}, blind readers expressed a mean level of agreement of $\mu=6.82$ ($SD=0.38$), indicating overwhelming support for visualization descriptions. 
All blind readers reported encountering inaccessible visualizations: either multiple times a week ({43}\%, n=13), everyday ({20}\%, n=6), once or twice a month ({20}\%, n=6), or at most once a week ({17}\%, n=5).
These readers reported primarily encountering these barriers on social media ({30}\%, n=9), on newspaper websites ({13}\%, n=4), and in educational materials ({53}\%, n=16)\,---\,but, most often, barriers were encountered in all of the above contexts ({53}\%, n=16).
% Readers report encountering inaccessible visualizations in each of the following contexts: {30}\% (n=9) on social media (Twitter, Facebook), {13}\% (n=4) on newspaper websites, {3}\% (n=1) in educational materials, and {53}\% (n=16) in all of the above.
Blind readers overwhelmingly agreed with the statements 
\say{I often feel that important public information is inaccessible to me, because it is only available in a visual format} ($\mu=6.1$, $\sigma=1.49$), and \say{Providing textual descriptions of data visualizations is important to me} ($\mu=6.83$, $\sigma=0.38$). 

\begin{displayquote}
% \small
    \vspace{-1mm}
   \hquoteA{I am totally blind, and virtually all data visualizations I encounter are undescribed, and as such are unavailable. This has been acutely made clear on Twitter and in newspapers around the COVID-19 pandemic and the recent U.S. election.
   Often, visualizations are presented with very little introduction or coinciding text. I feel very left out of the world and left out of the ability to confidently traverse that world. The more data I am unable to access, the more vulnerable and devalued I feel.} (BR5)
    \vspace{-1mm}
\end{displayquote}

\noindent
By contrast, sighted readers neither agreed nor disagreed regarding the inaccessibility of information conveyed visually ($\mu=4$, $\sigma=1.57$).
Similarly, they were split on whether they ever experienced barriers to reading visualizations, with {52}\% (n={47}) reporting that they sometimes do (especially when engaging with a \emph{new} topic) and {48}\% (n={43}) reporting that they usually do not.
Nevertheless, sighted readers expressed support for natural language descriptions of visualizations ($\mu=5.60$, $\sigma=1.27$).
A possible explanation for this support is that\,---\,regardless of whether the visualization is difficult to read\,---\,descriptions can still facilitate comprehension. For instance, SR64 noted that \hquoteB{textual description requires far less brainpower and can break down a seemingly complex visualization into an easy to grasp overview.}

% ``I've worked on different data science projects and studied statistics in college. I have the hardest time reading data visualizations when they have a tiny format or do not explain what each axis represents.''

%----------------------------------------------------------------
% \textbf{(B)~\emph{Reader Groups Disagree About Contextual Content.}}
\subsubsection{Reader Groups Disagree About Contextual Content}
\label{section:qualResultsB}
%----------------------------------------------------------------
A majority of blind readers ({63}\%, n=19) were emphatic that descriptions should \emph{not} contain an author's subjective interpretations, contextual information, or editorializing about the visualized data (i.e., Level 4 content).
Consistent with blind readers ranking this as among the least useful (Region~\ref{table:evaluation}.B), BR20 succinctly articulated a common sentiment: \hquoteA{I want the information to be simply laid out, not peppered with subjective commentary... I just prefer it to be straight facts, not presumptions or guesstimates.} BR4 also noted that an author's \hquoteA{opinions} about the data \hquoteA{should absolutely be avoided,} and BR14 emphasized {agency} when interpreting data: \hquoteA{I want to have the time and space to interpret the numbers for myself before I read the analysis.}
% \begin{displayquote}
%   \emph{\say{I don't want a description to be analytical and to include text about the possible causes of what the visualization shows. That should be separate. I want to have the time and space to interpret the numbers for myself before I read the analysis.}} (BP14)
% \end{displayquote}
% \begin{displayquote}
%   \emph{\say{I want the information to be simply laid out, not peppered with subjective commentary that may already be found in the article. While the information is useful and ties into that visualization, I just prefer it to be straight facts, not presumptions or guesstimates.
%   %I've found the most informative is laying out the numbers in a well thought out manner without being too overwhelming.
%   }} (BP20)
% \end{displayquote}
By contrast, many sighted readers {41}\% (n={37}) expressed the opposite sentiment (Region~\ref{table:evaluation}.F) noting that, for them, the most useful descriptions often \hquoteB{told a story,} communicated an important conclusion, or provided deeper insights into the visualized data.
As SR64 noted: \hquoteB{A description that simply describes the visualization and its details is hardly useful, but a description that tells a story using the data and derives a solution from it is extremely useful.}
Only {4}\% (n={4}) of sighted readers explicitly stated that a description should exclude Level 4 semantic content.

% \begin{displayquote}
%   \emph{\say{A description that simply describes the visualization and its details is hardly useful, but a description that tells a story using the data and derives a solution from it is extremely useful.
%   }} (SR64)
% \end{displayquote}

%----------------------------------------------------------------
\subsubsection{Some Readers Prefer Non-Statistical Content}
\label{section:qualResultsC}
%----------------------------------------------------------------
Overall, blind readers consistently ranked both Levels 2 and 3 as the most useful (Region~\ref{table:evaluation}.C). But, some readers explicitly expressed preference for the latter over the former, highlighting two distinguishing characteristics of Level 3 content: that it conveys not only descriptive statistics but overall perceptible trends, and that it is articulated in commonplace or ``natural''-sounding language.
For instance, BR26 remarked that a visualization description is \hquoteA{more useful if it contains the summary of the overall trends and distributions of the data rather than just mentioning some of the extreme values or means.}
Similarly, BR21 noted that \hquoteA{not everyone who encounters a data visualization needs it for statistical purposes,} and further exclaimed \hquoteA{I want to know how a layperson sees it, not a statistician; I identify more with simpler terminology.}
% Regarding the second characteristic, BR21 said the following to say.
% \begin{displayquote}
%   \emph{\say{Not everyone who encounters a data visualization needs it for statistical purposes... I want to know how a layperson sees it, not a statistician. I identify more with simpler terminology.}} (BR21)
% \end{displayquote}
These preferences help to further delineate Level 3 from Levels 2 and 4.
Content at Level 3 is ``non-statistical'' in the sense that it does only report statistical concepts and relations (as in Level 2), but neither does it do away with statistical ``objectivity'' entirely, so as to include subjective interpretation or speculation (as content in Level 4 might).
In short, Level 3 content conveys statistically-grounded concepts in not-purely-statistical terms, a challenge that is core to visualization, and science communication more broadly.

% \begin{displayquote}
%   \emph{\say{
%     I think a textual description is more useful if it contains the summary of the overall trends and distributions of the data rather than just mentioning some of the extreme values or means.}} (BR26)
% \end{displayquote}

%----------------------------------------------------------------
\subsubsection{Combinations of Content Levels Are Likely Most Useful}
\label{section:qualResultsD}
%----------------------------------------------------------------
While roughly {12}\% readers from both blind and sighted groups indicated that a description should be as concise as possible, among blind readers, {40}\% (n={12}) noted that the most useful descriptions would combine content from multiple levels.
% "shorter is better": blind n=4/30, sighted n=11/90
This finding helps to explain the bimodality in Level 1 rankings we identified in the previous section. 
According to BR9, Level 1 content is only useful if other information is also conveyed: \hquoteA{All of the descriptions provided in this survey which *only* elaborated on x/y and color-coding are almost useless.}
% \begin{displayquote}
%   \emph{\say{Descriptions of what is represented on the axes is only useful if other information is conveyed. All of the descriptions provided in this survey which *only* elaborated on x/y and color-coding are almost useless.
%   }} (BR9)
% \end{displayquote}
This sentiment was echoed by BR5, who added that if Level 1 content were \hquoteA{combined with the [Level 2 or Level 3], that'd make for a great description.}
% \begin{displayquote}
%   \emph{\say{I’d liked most of the descriptions, actually, if they were ``combined.'' For instance, some of the descriptions (I mostly rated as [least useful]) only mentioned the title and labels of the axes. If combined with the [Levels 2 or Level 3], that'd make for a great description.
%   }} (BR5)
% \end{displayquote}
% \noindent
This finding has implications for research on automatic visualization captioning: these methods should aim to generate not only the lower levels of semantic content, but to more richly communicate a chart’s overall trends and statistics, sensitive to reader preferences.

\subsubsection{Some Automatic Methods Raise Ethical Concerns}
\label{section:qualResultsE}
Research on automatically generating visualization captions is often motivated by the goal of improving information access for people with visual disabilities~\cite{qian_generating_2021, demir_summarizing_2012, obeid_chart--text_2020, qian_formative_2020}.
However, when deployed in real-world contexts, these methods may not confer their intended benefits, as one blind reader in our evaluation commented.
\begin{displayquote}
% \small
    \vspace{-1mm}
    \hquoteA{A.I. attempting to convert these images is still in its infancy. Facebook and Apple auto-descriptions of general images are more of a timewaster than useful. As a practical matter, if I find an inaccessible chart or graph, I just move on.} (BP22)
    \vspace{-1mm}
\end{displayquote}
Similarly, another participant (BR26) noted that if a description were to only describe a visualization's encodings then \hquoteA{the reader wouldn't get any insight from these texts, which not only increases the readers' reading burden but also conveys no effective information about the data.}
% \begin{displayquote}
%   \emph{\say{
%     The description would be almost useless if it only describes the encoding of the visualization because the reader wouldn't get any insight from these texts, which not only increases the readers' reading burden but also conveys no effective information about the data.}} (BR26)
% \end{displayquote}
These sentiments reflect some of the ethical concerns surrounding the deployment of nascent \ac{CV} and \ac{NLP} models, which can output accurate but minimally informative content\,---\,or worse, can output erroneous content to a trusting audience~\cite{macleod_understanding_2017, obeid_chart--text_2020}.
Facebook's automatic image descriptions, for example, have been characterized by technology educator Chancey Fleet as \say{\emph{famously useless in the Blind community}} while \say{\emph{garner[ing] a ton of glowing reviews from mainstream outlets without being of much use to disabled people}}~\cite{fleet_things_2021, hanley_computer_2021}. Such concerns might be mitigated by developing and evaluating automatic methods with disabled readers, through participatory design processes~\cite{lundgard_sociotechnical_2019}.

% Chancey Fleet. "Things which garner a ton of glowing reviews from mainstream outlets without being of much use to disabled people. for instance, facebook’s auto image descriptions, much loved by sighted journos but famously useless in the blind community," January 2021.